\documentclass[11pt,a4paper]{article}

% ── Page geometry ──
\usepackage[hmargin=0.9in,vmargin=1in]{geometry}

% ── Language ──
\usepackage[english]{babel}

% ── Microtypography (better justification, fewer overfull boxes) ──
\usepackage{microtype}

% ── Modern fonts via fontspec (XeLaTeX) ──
\usepackage{fontspec}
\setmainfont{Roboto}[
  Path          = /home/devakchaudhari/.fonts/,
  Extension     = .ttf,
  UprightFont   = *-Regular,
  BoldFont      = *-Medium,
  ItalicFont    = *-Regular,
  ItalicFeatures = {FakeSlant=0.167},
]
\setmonofont{Liberation Mono}[
  Scale=MatchLowercase,
]

% ── Mathematics ──
\usepackage{amsmath,amssymb,amsthm}
\usepackage{unicode-math}
\setmathfont{STIX Two Math}[
  Scale=MatchLowercase,
]

% ── Graphics & floats ──
\usepackage{graphicx}
\usepackage{float}
\usepackage{caption}
\usepackage{subcaption}

% ── Code listings ──
\usepackage{listings}
\usepackage{xcolor}
\definecolor{codegreen}{rgb}{0,0.4,0.2}
\definecolor{codegray}{rgb}{0.5,0.5,0.5}
\definecolor{codecyan}{rgb}{0,0.5,0.5}
\definecolor{codeyellow}{rgb}{0.55,0.5,0.1}
\lstset{
  basicstyle={\footnotesize\ttfamily},
  commentstyle=\color{codegreen},
  keywordstyle=\color{codecyan},
  stringstyle=\color{codeyellow},
  numbers=left,
  numberstyle=\tiny\color{codegray},
  frame=single,
  breaklines=true,
}

% ── Hyperlinks ──
\usepackage[hidelinks]{hyperref}

% ── Bibliography ──
\usepackage[backend=bibtex,sorting=none]{biblatex}
\addbibresource{references.bib}

% ── Custom commands ──
\newcommand{\trinary}{\texttt{Trinary}}

\title{Trinary: A Ternary (Base-3) Computer Simulation in Pure Python}
\author{}
\date{\today}



\begin{document}

\maketitle
\thispagestyle{empty}

\begin{abstract}
\trinary{} is a complete ternary (base-3) computer system simulated entirely in
pure Python. The project spans the full abstraction hierarchy --- from
fundamental logic gates (TNOT, TAND, TOR) and a ripple-carry adder, through a
general-purpose CPU with 26 instructions, memory-mapped I/O, interrupts, and a
hardware timer, up to a minimal operating system shell and a PyQt6 desktop
visual simulator. This document describes the architecture, instruction set,
memory model, interrupt system, and implementation decisions behind the
simulator, and serves as a reference for further research into non-binary
computing models.
\end{abstract}

\tableofcontents
\clearpage

% ===================================================================
\section{Introduction to Ternary Computing}
\label{sec:introduction}
% ===================================================================

\subsection{Binary vs.\ Ternary Number Systems}

Modern computing is overwhelmingly binary.  Since the earliest electronic
computers --- ENIAC (1945), EDVAC, and the von Neumann architecture --- the
two-state switch (on/off, true/false, 1/0) has dominated digital logic.
Claude Shannon's 1937 master's thesis~\cite{shannon1937relay} established
that Boolean algebra could be realized with electrical switching circuits,
cementing binary as the natural choice for hardware.

Yet binary is not the only possibility.  A ternary (base-3) digit --- called
a \emph{trit} --- can represent three distinct states, typically denoted
0, 1, and 2.  This extra state per digit has profound implications for
information density, circuit complexity, and arithmetic efficiency.

\begin{table}[H]
  \centering
  \caption{Comparing the number of digits needed to represent values in
           binary vs.\ ternary.}
  \label{tab:digit-comparison}
  \begin{tabular}{c|c|c|c}
    \textbf{Value} & \textbf{Binary digits} & \textbf{Ternary digits}
    & \textbf{Savings} \\ \hline
    10    & 4  & 3  & 25\% \\
    100   & 7  & 5  & 29\% \\
    1000  & 10 & 7  & 30\% \\
    10$^6$ & 20 & 13 & 35\% \\
    10$^8$ & 27 & 17 & 37\%
  \end{tabular}
\end{table}

\subsection{Information-Theoretic Advantages}

Donald Knuth, in \emph{The Art of Computer Programming}, notes that the
base $r$ that maximizes information density is $e$, with 3 being the
closest integer~\cite{knuth2003balanced}.  The \emph{radix economy}
$E(r, N) = r \cdot \lceil \log_r N \rceil$ measures the cost of
representing a number $N$ in base $r$.  Base 3 minimizes this cost among
integer bases:

\begin{itemize}
  \item Binary $(r=2)$: $2 \cdot \log_2 N$
  \item Ternary $(r=3)$: $3 \cdot \log_3 N \approx 3 \cdot \frac{\ln N}{\ln 3}
        \approx 2.73 \cdot \ln N$
  \item Quaternary $(r=4)$: $4 \cdot \log_4 N \approx 4 \cdot
        \frac{\ln N}{\ln 4} \approx 2.89 \cdot \ln N$
\end{itemize}

Ternary achieves the lowest constant factor ($\approx 2.73$) among integer
bases, meaning it can represent a given range of values with the fewest
physical components in an idealized circuit model.  Hurst further explores
these tradeoffs in the context of multiple-valued logic
systems~\cite{hurst1990multiple}.

\subsection{The Soviet Setun Computer (1958)}

The most famous practical implementation of ternary computing is the
\texttt{Setun} computer, developed at Moscow State University by
Nikolai~P.~Brusentsov and his team in 1958~\cite{brusentsov1965setun}.

The Setun used \emph{balanced ternary} (trits valued $-1, 0, +1$), which
offered a natural representation of negative numbers without a dedicated
sign bit.  Key specifications included:

\begin{itemize}
  \item 36-bit (trit) word length --- approximately 18 decimal digits of
        precision
  \item Ferrite-diode memory cells for the register file
  \item Magnetic drum memory for bulk storage
  \item Approximately 50 machines were produced and deployed in Soviet
        universities and research laboratories
\end{itemize}

The Setun demonstrated that ternary computing was not merely a theoretical
curiosity but a viable engineering alternative.  Its production ended in
the mid-1960s as binary hardware --- simpler to design with
transistor-transistor logic (TTL) --- became dominant.

\subsection{Why Simulate a Ternary Computer in Software?}

Building a physical ternary computer today would require custom logic
families (e.g., CNTFET-based ternary gates~\cite{lin2022ternary}) and
fabrication processes that remain niche.  However, a software simulation
offers several advantages:

\begin{enumerate}
  \item \textbf{Pedagogy}: A simulated machine can be inspected at every
        level --- from individual trits and logic gates up to running
        programs and an operating system --- making it an ideal teaching
        tool for computer architecture.
  \item \textbf{Experimentation}: Instruction sets, memory models,
        addressing modes, and interrupt systems can be prototyped,
        modified, and benchmarked without hardware constraints.
  \item \textbf{Reproducibility}: A pure-Python simulator runs on any
        platform with a Python interpreter (3.10+), ensuring that results
        are fully deterministic and shareable.
  \item \textbf{Accessibility}: No specialized hardware or FPGA toolchains
        required --- a single \texttt{pip install} is sufficient to boot
        the simulated machine and interact with its OS shell.
\end{enumerate}

\subsection{Scope of This Book}

This book traces the complete stack of the \trinary{} ternary computer
simulation, progressing through 25 chapters from first principles to a
working graphical debugger:

\begin{enumerate}
  \item[(Ch.\ 2)] System architecture and module dependency graph
  \item[(Ch.\ 3)] Environment setup and ``Hello, Ternary''
  \item[(Ch.\ 4--5)] The Trit datatype and base-conversion systems
  \item[(Ch.\ 6--7)] Ternary logic gates and adder circuits
  \item[(Ch.\ 8--9)] Arithmetic via decimal round-trip and the ALU
  \item[(Ch.\ 10--11)] Register file and memory architecture
  \item[(Ch.\ 12--20)] CPU design, all 26 instructions, interrupts, and timer
  \item[(Ch.\ 21--22)] Two-pass assembler and machine-code encoding
  \item[(Ch.\ 23)] Memory-mapped display and keyboard subsystem
  \item[(Ch.\ 24)] The operating system shell
  \item[(Ch.\ 25)] The PyQt6 visual simulator
\end{enumerate}

Each chapter includes runnable code examples that readers can execute
immediately against the \trinary{} package.

\subsection{Verifying Your Environment}

Before proceeding, verify that your Python environment is ready:

\begin{lstlisting}[language=bash, caption={Installing \trinary{} and running
the test suite.}]
pip install -e .
python -m pytest tests/ -v
\end{lstlisting}

The test suite (113 tests, all passing) exercises every component from
base conversion to interrupt handling.  A successful run produces output
similar to:

\begin{lstlisting}[basicstyle=\ttfamily\footnotesize]
======================== 113 passed in 0.06s =========================
\end{lstlisting}

You can also launch the interactive OS shell immediately:

\begin{lstlisting}[language=bash]
python -m trinary.os
\end{lstlisting}

If the shell boots and displays a \texttt{Trishell!>} prompt, your
environment is fully functional and you are ready to begin.

% ===================================================================
\section{Data Representation}
\label{sec:data}
% ===================================================================

All numeric values are represented as signed-magnitude ternary strings.  A
leading minus sign (\texttt{-}) denotes a negative number; its absence denotes
a non-negative number.  The \texttt{Trit} class encapsulates a single ternary
digit with possible values 0, 1, 2.

\begin{table}[H]
  \centering
  \caption{Ternary-decimal correspondence.}
  \label{tab:ternary-decimal}
  \begin{tabular}{c|c}
    \textbf{Ternary} & \textbf{Decimal} \\ \hline
    0      & 0  \\
    1      & 1  \\
    2      & 2  \\
    10     & 3  \\
    11     & 4  \\
    12     & 5  \\
    20     & 6  \\
    21     & 7  \\
    22     & 8  \\
    100    & 9  \\
    -10    & $-3$ \\
    -101   & $-10$
  \end{tabular}
\end{table}

Conversion functions provide bidirectional translation between ternary strings,
decimal integers, and binary integers, enabling seamless interaction with
Python's native numeric operations.

% ===================================================================
\section{Logic Gates}
\label{sec:gates}
% ===================================================================

Ternary logic extends Boolean logic by introducing a third truth value.  Three
primitive gates are defined:

\begin{align*}
  \text{TNOT}(x) &= 2 - x \\
  \text{TAND}(x, y) &= \min(x, y) \\
  \text{TOR}(x, y) &= \max(x, y)
\end{align*}

where $x, y \in \{0, 1, 2\}$.  Truth tables for each gate are available
through the \texttt{trinary.logic} module.

% ===================================================================
\section{Addition Paths}
\label{sec:addition}
% ===================================================================

The project provides two parallel addition implementations:

\paragraph{Native adder (\texttt{adder.py})}
A half-adder, full-adder, and ripple-carry adder operating directly on ternary
digits (values 0, 1, 2).  The ripple-carry adder chains full-adders for each
trit position and does not handle negative numbers.

\paragraph{Decimal round-trip (\texttt{arithmetic.py})}
The ADD, SUB, MUL, and DIV functions convert ternary operands to decimal,
perform the operation in Python integers, and convert the result back to
ternary.  This path naturally handles signed-magnitude negatives.

The ALU (\texttt{alu.py}) selects the latter path for all arithmetic
operations.

% ===================================================================
\section{CPU Architecture}
\label{sec:cpu}
% ===================================================================

\subsection{Registers}

The CPU provides four general-purpose registers and three special-purpose
registers:

\begin{table}[H]
  \centering
  \caption{Register file.}
  \label{tab:registers}
  \begin{tabular}{c|l}
    \textbf{Name} & \textbf{Purpose} \\ \hline
    R0   & Primary accumulator \\
    R1   & Secondary register \\
    R2   & General purpose \\
    R3   & General purpose (OS convention: always 0) \\
    PC   & Program counter \\
    SP   & Stack pointer (starts at 255, grows downward) \\
    Flags & Condition flags: ZERO, EQUAL, GREATER, LESS
  \end{tabular}
\end{table}

\subsection{Memory Map}

The address space is 512 trits wide.  Memory is divided into regions:

\begin{table}[H]
  \centering
  \caption{Memory map.}
  \label{tab:memory}
  \begin{tabular}{c|l}
    \textbf{Addresses} & \textbf{Purpose} \\ \hline
    0--127   & General data / program \\
    128--199 & Stack \\
    200--255 & Video RAM (56 characters, $7 \times 8$ grid) \\
    260      & Keyboard buffer \\
    256--511 & Reserved / future expansion
  \end{tabular}
\end{table}

\subsection{Fetch-Decode-Execute Cycle}

\subsubsection{Fetch}
The instruction at \texttt{program[PC]} is read as a raw string.

\subsubsection{Decode}
The instruction string is split into tokens.  The first token is the opcode;
subsequent tokens are operands (register names, numeric literals, or labels).

\subsubsection{Execute}
The CPU dispatches to the appropriate handler for the decoded opcode.  Side
effects include register changes, memory writes, stack operations, and PC
modification.

\subsubsection{Update}
Unless the instruction modified \texttt{PC} (e.g., JMP, JZ, CALL), \texttt{PC}
is incremented by one.

\subsection{Dual Stacks}

The CPU maintains two distinct stacks:

\begin{description}
  \item[Memory stack] (SP, addresses 128--199) --- used by PUSH and POP.
        Grows downward from address 255.
  \item[Call stack] (Python list) --- used by CALL and RET for subroutine
        return addresses.  Not visible in the memory address space.
\end{description}

This design separates data stacking from control-flow stacking, simplifying
implementation and avoiding address-space conflicts.

\subsection{Flags}

The CMP instruction compares two ternary values and sets four boolean flags:
\begin{itemize}
  \item \texttt{ZERO}: result is zero.
  \item \texttt{EQUAL}: operands are equal.
  \item \texttt{LESS}: first operand $<$ second.
  \item \texttt{GREATER}: first operand $>$ second.
\end{itemize}
Jump instructions (JZ, JNZ) and conditional logic use these flags.

\subsection{Interrupts}

The interrupt system provides:

\begin{itemize}
  \item \textbf{Interrupt vector table (IVT)}: 8 entries, each storing a
        handler address.  Set via SETIVT.
  \item \textbf{Software interrupts}: triggered by the INT instruction with an
        interrupt number argument.
  \item \textbf{Hardware timer}: periodic interrupt (interrupt 0) configured
        via SETTIMER.  Can be enabled/disabled with EI/DI.
  \item \textbf{Interrupt return}: IRET pops the saved PC and re-enables
        interrupts.
\end{itemize}

% ===================================================================
\section{Instruction Set}
\label{sec:isa}
% ===================================================================

The CPU implements 26 instructions, grouped into functional categories:

\subsection{Data Movement}
LOAD, MOV, CLR

\subsection{Arithmetic}
ADD, SUB, MUL, DIV

\subsection{Logical}
AND, OR, NOT

\subsection{Comparison}
CMP

\subsection{Control Flow}
JMP, JZ, JNZ

\subsection{Subroutines}
CALL, RET

\subsection{Stack}
PUSH, POP

\subsection{Memory}
STOREM, LOADM

\subsection{Interrupts}
INT, IRET, EI, DI, SETIVT

\subsection{Timer}
SETTIMER

\subsection{System}
HALT

See \texttt{docs/instruction-set.md} for the complete opcode reference.

% ===================================================================
\section{Assembler}
\label{sec:assembler}
% ===================================================================

The assembler (\texttt{assembler.py}) is a two-pass assembler:

\begin{enumerate}
  \item \textbf{Pass 1}: Scan all source lines for label definitions
        (e.g., \texttt{start:}).  Record the address (line number) of each
        label in a symbol table.
  \item \textbf{Pass 2}: Re-scan all source lines, resolve label references
        in jump/call operands using the symbol table, and emit a list of
        instruction strings.
\end{enumerate}

Inline comments are supported with both \texttt{\#} and \texttt{;} delimiters,
handled by \texttt{assembler.parse\_line}.

% ===================================================================
\section{Machine Code}
\label{sec:machine-code}
% ===================================================================

The \texttt{machine.py} module provides a machine-code encoder and decoder that
converts between assembly strings and variable-length ternary opcode strings.
Currently, 21 of the 26 instructions have machine-code mappings.  The
instructions INT, IRET, EI, DI, SETIVT, and SETTIMER do not yet have opcode
assignments in the encoding table.

% ===================================================================
\section{OS Shell}
\label{sec:os}
% ===================================================================

The minimal operating system (\texttt{os.py}) is written entirely in assembly
(~342 lines).  It provides an interactive command shell with the following
commands:

\begin{table}[H]
  \centering
  \caption{OS shell commands.}
  \label{tab:os-commands}
  \begin{tabular}{l|l}
    \textbf{Command} & \textbf{Description} \\ \hline
    HELP     & Lists available commands \\
    CLS      & Clears the display \\
    REGS     & Shows register values \\
    MEMDUMP  & Dumps first 8 memory addresses \\
    RUN      & Placeholder: prints ``RUN!''
  \end{tabular}
\end{table}

The OS implements memory-mapped keyboard polling, an input line buffer with
backspace support, cursor management with display scrolling, and
multi-character command parsing.

% ===================================================================
\section{Display System}
\label{sec:display}
% ===================================================================

The display system maps memory addresses 200--255 to a 7-row by 8-column
character display (56 characters total).  A \texttt{PixelDisplay} class
provides a 27$\times$27 pixel framebuffer representation.  Address 260 serves
as a keyboard buffer: the OS polls this address, and the UI (or a human user)
writes ASCII key codes into it.

% ===================================================================
\section{PyQt6 Desktop UI}
\label{sec:ui}
% ===================================================================

The \texttt{trinary.ui} package provides a PyQt6-based graphical simulator
with the following panels:

\begin{itemize}
  \item \textbf{Assembler editor}: syntax highlighting (opcodes=cyan,
        registers=yellow, labels=purple, comments=green, numbers=white),
        clickable breakpoints in the gutter.
  \item \textbf{Machine code viewer}: ternary opcode strings displayed
        alongside assembly source.
  \item \textbf{Register panel}: R0--R3 values with flash animation on change;
        PC, SP, and four flag indicators.
  \item \textbf{Memory viewer}: 512-row table with colour-coded access tracking
        (green=write, blue=read, current PC highlight).
  \item \textbf{Stack viewer}: live stack contents with SP position marker.
  \item \textbf{Execution trace}: step-by-step history (step number, PC,
        instruction, registers, flags).
  \item \textbf{Pixel display}: 27$\times$27 pixel framebuffer with keyboard
        capture.
  \item \textbf{Toolbar}: demo selector (13 programs) and execution controls
        (Assemble, Run, Step Into, Step Over, Reset, Pause, Continue).
\end{itemize}

The UI is fully decoupled from the back-end --- it interacts exclusively
through the public CPU and memory APIs.  The stylesheet
(\texttt{styles.py}) provides a dark cyberpunk theme.

% ===================================================================
\section{Testing}
\label{sec:tests}
% ===================================================================

The project includes 113 tests across 8 test files:

\begin{table}[H]
  \centering
  \caption{Test suite.}
  \label{tab:tests}
  \begin{tabular}{l|l}
    \textbf{File} & \textbf{Coverage} \\ \hline
    test\_conversion.py  & Trit class, base conversion \\
    test\_arithmetic.py  & ADD, SUB, MUL, DIV \\
    test\_alu.py         & All 8 ALU operations \\
    test\_assembler.py   & Label resolution, comments \\
    test\_cpu.py         & All 26 opcodes, flags, stack, interrupts \\
    test\_cpu\_stress.py & Nested calls, long loops \\
    test\_display.py     & Memory-mapped display, keyboard
  \end{tabular}
\end{table}

Tests can be run with:
\begin{lstlisting}[language=bash]
python -m pytest tests/ -v
\end{lstlisting}

% ===================================================================
\section{Related Work}
\label{sec:related}
% ===================================================================

The most famous ternary computer is the Soviet \texttt{Setun}
(1958)~\cite{brusentsov1965setun}, which used balanced ternary
($-1, 0, +1$) and was built from ferrite-diode cells.  More recently,
ternary logic has been explored in the context of
\textsc{cntfet}-based circuits~\cite{lin2022ternary} and
memristive ternary systems.

Software simulators of non-binary architectures include the
\texttt{ternary-computer} project and various
university-level teaching simulators.  \trinary{} distinguishes itself by
combining a complete, documented ISA with an interactive OS shell and a
modern desktop UI, all implemented in a single, well-tested Python
package.

% ===================================================================
\section{Future Directions}
\label{sec:future}
% ===================================================================

Possible extensions and improvements:

\begin{itemize}
  \item Complete machine-code mappings for the remaining six instructions
        (INT, IRET, EI, DI, SETIVT, SETTIMER).
  \item A balanced ternary mode ($-1, 0, +1$) as an alternative to
        signed-magnitude representation.
  \item An expanded memory model supporting paging or banking.
  \item A network interface or I/O port system.
  \item A compiler from a high-level language (e.g., a subset of C) to
        ternary assembly.
  \item WebAssembly compilation via Pyodide for in-browser simulation.
\end{itemize}

% ===================================================================
\printbibliography
% ===================================================================

\end{document}
